\subsection{Effective countable concordance models}

Classical decision procedures over real closed fields and the Moore
description of simplicial-group homotopy give an effective version of
the finite-model comparisons.  The output is a complete countable
presentation; it need not solve the word problem.

\begin{theorem}\label{thm:effective-countable-concordance-models}
Given an ordered finite simplex list for a finite abstract complex $K$
whose realization is a closed PL manifold, one can construct countable
simplicial groups $G^{\mathrm{PL}}(K)$ and $G^{\mathrm{SA}}(K)$ with
the following properties.  Their simplices have finite integer codes;
validity, equality and degeneracy are decidable.  All group operations,
faces and degeneracies are computable by terminating algorithms, as
is a simplicial homomorphism
\[
                   G^{\mathrm{PL}}(K)\longrightarrow G^{\mathrm{SA}}(K).
\]
For every real closed field $R$ the groups, interpreted over the real
algebraic numbers, give weak equivalences
\[
 \begin{aligned}
 G^{\mathrm{PL}}(K)&\longrightarrow P_{\mathrm{PL},R}(|K|_R),\\
 G^{\mathrm{SA}}(K)&\longrightarrow P_{\mathrm{SA},R}(|K|_R).
 \end{aligned}
\]
In these concordance groups the total map preserves its entire
parameter simplex, fixes the bottom pointwise, and preserves the
free top setwise. In the PL case the total map and its inverse
are finite piecewise linear; in the SA case they are semialgebraic.
In both cases they are jointly continuous on the whole product.

Uniformly in $K$ and $r\geq0$, an algorithm enumerates a complete
countable group presentation of each of these homotopy groups.
For $r\geq1$ it can instead enumerate a countable integer relation
matrix for the resulting abelian group, based at the identity.
Translation gives the homotopy groups at other components.
The PL presentation in degree zero also presents the stipulated
definable component group in every fixed o-minimal expansion.
If every component of $|K|$ has dimension at least five, the PL
construction is also a weak model for the real TOP concordance
space and presents all of its homotopy groups.

The empty complex gives the constant one-element simplicial group.
The manifold condition is a promise on the input; no recognition
algorithm for closed PL manifolds is asserted.  The presentations
need not be finite, and no terminating word problem, invariant-factor
calculation or complexity bound is claimed.
\end{theorem}

\begin{proof}
Let $F=\mathbb Q^{\mathrm{rc}}$.  Realize the ordered vertices of $K$
as distinct standard basis vectors, and put
\[
 B=|K|_F,\qquad Z=B\times[0,1],\qquad
 A=B\times\{0\},\qquad T=B\times\{1\}.
\]
These are finite rational polyhedra.  A $q$-simplex will be described
by a graph in the variables $(z,z',t)$, with $z,z'\in Z$ and
$t\in\Delta^q_F$.  The repeated parameter is implicit: its total
map is $(z,t)\mapsto(z',t)$.

\medskip\noindent\textit{Effective scalar and formula operations.}
Represent an element of $F$ by a nonzero integer polynomial and a
rational open interval containing exactly one of its real roots,
with neither endpoint a root.  Such codes are enumerable and their
validity is decidable by Sturm's theorem.  A finite list of algebraic
constants can alternatively be replaced in any formula by variables
subject to their root and isolating-interval conditions.  These
conditions have unique solutions.  Decision and quantifier
elimination over the rational coefficients therefore decide equality,
order and all first-order tests with the encoded constants.
To compute a scalar sum, product or inverse of a nonzero scalar,
enumerate valid scalar codes and test the appropriate equality; a
code exists and the search terminates.  These are classical effective
procedures; the exact locators used here are
\cite[online revision of 26 June 2016, Theorem~2.62, p.~62;
Algorithm~14.16, pp.~586--587; Algorithm~14.21 and Theorem~14.22,
pp.~591--593]{BasuPollackRoy2016Online}.

We also use the elementary linear version of elimination.  Put a
Boolean combination of affine equalities and inequalities in
disjunctive normal form.  Equalities with a nonzero eliminated
coefficient allow substitution.  In the other cases the inequalities
give finitely many upper and lower bounds for that variable.  Existence
means that every lower bound is at most every upper bound, strictly
so whenever either bound was strict.  With no bound on one side
there is no additional condition.  This eliminates one existential
variable; negation handles universal quantifiers.  Repetition gives
a terminating linear decision and elimination procedure over $F$.

\medskip\noindent\textit{PL graphs and their validity.}
A raw PL code is a finite list of closed simplices in the graph
coordinates, with vertices in $F$ and with affine independence
checked.  It describes the union $\Gamma$ of these simplices.
Barycentric coordinates express membership by a linear formula.
Require the following five decidable properties:
\begin{enumerate}
\item $\Gamma\subset Z\times Z\times\Delta^q_F$;
\item for each $(z,t)\in Z\times\Delta^q_F$ there is exactly one
      $z'$ with $(z,z',t)\in\Gamma$;
\item for each $(z',t)\in Z\times\Delta^q_F$ there is exactly one
      $z$ with $(z,z',t)\in\Gamma$;
\item $(z,z,t)\in\Gamma$ for every $z\in A$ and $t\in\Delta^q_F$;
\item on $\Gamma$, $z\in T$ if and only if $z'\in T$.
\end{enumerate}
All these statements are linear first-order sentences.  Equality
of two coded graphs is a linear sentence as well.

The two graph projections of a valid code are PL homeomorphisms.
Indeed, each is injective on each graph simplex, hence on its affine
hull: a nonzero kernel direction there would give two distinct points
in the relative interior with the same image.  Its inverse on the
projected simplex is therefore affine.  The projected simplices are
closed, form a finite cover, and their inverses agree on intersections.
Finite closed pasting gives continuity of both total maps over $F$.
The finite polyhedral family has a compatible triangulation over $F$:
intersect the finitely many face hyperplanes, then triangulate the
resulting polyhedral cells compatibly, for instance by their face
barycentres.  Thus the map and its inverse are jointly PL.  This
argument uses no ordinary compactness assertion for the ordered field.
Conversely, every jointly PL concordance on the total product has a
finite graph triangulation of this kind, so the raw codes represent
exactly $P_{\mathrm{PL},F}(B)_q$.

Encode the finite strings in a fixed finite alphabet and order them
first by length and then lexicographically.  Replace a valid raw code
by the first valid code in this order defining its graph.  It is
enough to examine its finitely many predecessors.  Validity and
graph equality make this a terminating canonicalization algorithm;
the canonical codes themselves form a decidable set.  They give
$G^{\mathrm{PL}}(K)_q$.

For multiplication, form the linear graph formula
\[
  \exists w\;\bigl(\Gamma_2(z,w,t)\ \wedge\ \Gamma_1(w,z',t)\bigr).
\]
For inversion interchange $z$ and $z'$. For a simplicial operator
$\alpha:\Delta^p_F\to\Delta^q_F$, define the pullback graph by
\[
 \Gamma_{\alpha^*g}(z,z',t)=
 [\,z\in Z,\ z'\in Z,\ t\in\Delta^p_F\,]
 \ \wedge\ \Gamma_g(z,z',\alpha(t)).
\]
The displayed domain conditions are part of the formula. All graph
equalities are tested in the full ambient coordinate space; in
particular, no points with $t\notin\Delta^p_F$ are admitted after an
affine parameter substitution. The identity is the diagonal graph
on $Z\times\Delta^q_F$. These operations are effective on formulas.
They describe jointly PL homeomorphisms: finite PL compositions,
inverses and these affine parameter pullbacks admit finite compatible
polyhedral subdivisions.  Enumerate valid raw codes and compare their
graphs with the resulting formula until equality holds, then
canonicalize.  Existence of such a graph code proves termination.
The identity has its diagonal graph code.  Consequently all operations
are total on valid inputs; an invalid input can be recognized and
rejected first.  Because the canonical graphs are actual maps,
the group and simplicial identities hold exactly on their codes.

\medskip\noindent\textit{Semialgebraic graphs.}
A raw SA code is a finite Boolean formula in polynomial equalities
and inequalities over the encoded field $F$.  Impose the same five
conditions and, in addition, continuity of the two total maps.
For example, the forward condition is the first-order sentence
\[
\begin{split}
 \forall (z,z',t)\in\Gamma\ \forall\varepsilon>0\ \exists\delta>0\
 \ \forall(w,w',u)\in\Gamma:\qquad\qquad\qquad\quad\\
 \|(w,u)-(z,t)\|^2<\delta^2
 \ \Longrightarrow\
 \|(w',u)-(z',t)\|^2<\varepsilon^2.
\end{split}
\]
Interchanging the source and target variables gives the inverse
condition.  This tests continuity in the parameter as well as in
the fibre.  Polynomial quantifier elimination decides all these
sentences and equality of two graphs.  Hence the same finite-predecessor
canonicalization gives a decidable set $G^{\mathrm{SA}}(K)_q$.
Every semialgebraic concordance gives such a graph and conversely.

Composition, inversion and parameter substitution are the same graph
operations as above.  Quantifier elimination puts their formulas in
the permitted raw form, after which canonicalization is effective.
These give the group and simplicial operations.  A PL graph has an
effective semialgebraic formula, so canonicalizing that formula gives
the required computable simplicial homomorphism.  In fact the two
constructed simplicial groups are coded copies of the respective
concordance groups over $F$.  Their asserted weak comparisons with
other real closed fields are exactly Theorem~\ref{result-40} and
Corollary~\ref{result-52}.

Degeneracy is also decidable in either model.  A $q$-simplex $a$,
with $q>0$, is degenerate precisely when
\[
                   a=s_i d_i a\quad\text{for some }0\leq i<q.
\]
If $a=s_i b$, the identity $d_i s_i b=b$ gives the displayed test;
the converse is immediate.  A degeneracy composite lies in the
image of some single $s_i$, and there are only finitely many indices
to test.  No zero-simplex is degenerate.

\medskip\noindent\textit{Complete homotopy presentations.}
Let $G$ be either coded simplicial group, and let $1$ denote the
identity in each degree.  Define decidable sets
\[
\begin{aligned}
 Z_0&=G_0,\\
 Z_r&=\{a\in G_r:d_i a=1\text{ for }0\leq i\leq r\}
                    \quad(r\geq1),\\
 N_{r+1}&=\{u\in G_{r+1}:d_i u=1\text{ for }1\leq i\leq r+1\}
                    \quad(r\geq0).
\end{aligned}
\]
The simplicial identities show $d_0N_{r+1}\subset Z_r$.
Conjugating $u\in N_{r+1}$ by $s_0a$, for $a\in Z_r$, stays in
$N_{r+1}$ and sends its first face to $a(d_0u)a^{-1}$.
Thus the image is a normal subgroup of $Z_r$.

The classical Moore description gives
\[
                 \pi_r(G,1)=Z_r/d_0N_{r+1},\qquad r\geq0,
\]
where degree zero means the component group.  For clarity about
conventions, May uses the last face: see
\cite[Theorem~17.1, pp.~67--68; both propositions numbered~17.3,
pp.~68--69; Proposition~17.4, p.~69]{May1967Simplicial}.
In the first-face convention an elementary homotopy from a cycle
$y$ to a cycle $x$ has a simplex $v$ with $d_0v=x$, $d_1v=y$,
and all remaining faces equal to $1$.  Then
$u=v(s_0y)^{-1}$ lies in $N_{r+1}$ and $d_0u=xy^{-1}$.
Conversely such a $u$ gives $v=u(s_0y)$.  Simplicial groups are
Kan, so this is precisely the usual cycle homotopy relation.
Multiplication of cycles gives the homotopy-group operation;
May's Proposition~17.2, p.~68, records this identification.
Theorem~16.6, p.~65, identifies these with the homotopy groups
of the realization.

Now enumerate the presentation
\[
 \left\langle x_a\ (a\in Z_r)\ \middle|\
 x_1=1,\quad x_ax_b=x_{ab}\ (a,b\in Z_r),\quad
 x_{d_0u}=1\ (u\in N_{r+1})\right\rangle.
\]
All generator and relator indices use the decidable tests and
terminating operations already constructed. At stage $m$, examine all
strings and pairs of strings of index at most $m$, printing the
valid cycle generators, multiplication relations, and boundaries.
Each stage terminates and every required item eventually occurs.
This does not wait for a next valid code when only finitely many exist.
The multiplication-table relations present $Z_r$ itself: the
assignments $x_a\mapsto a$ and $a\mapsto x_a$ are inverse
homomorphisms.  The boundary relations then quotient by exactly
the normal subgroup $d_0N_{r+1}$.  This proves completeness of the
presentation, not merely a surjection onto the desired group.

For $r\geq1$ this quotient is abelian, including $r=1$; this is
the first Proposition~17.3 on May's p.~68.  Therefore the equivalent
abelian presentation has free generators $e_a$ and relation rows
\[
                   e_1,\qquad e_a+e_b-e_{ab},\qquad e_{d_0u}.
\]
Every row has finite support and the rows can be enumerated uniformly.
This gives the asserted countable integer relation matrix.  The
boundary image is recursively enumerable; deciding membership in it
has not been used or asserted.

The degree-zero definable comparison and the real PL--TOP weak
comparison in the stated dimension range follow from
Theorem~\ref{result-4}.  They use a chosen real realization, not a
map of nonreal points into a TOP space.  The construction preserves
components of a disconnected base because the bottom is fixed.
The comparisons therefore apply componentwise.  For an empty base
use the one-element group in every degree.  This completes the
construction and all the assertions.
\end{proof}

The theorem constructs a computable PL-to-SA map but does not use that
map as an additional higher-comparison theorem.  Its full-space
identifications are the separately established coefficient comparisons.
